%% OVERFORE-Subdocumento1-Anexos.tex — generado desde OVERFORE-Subdocumento1-Anexos.docx
%% Formato: clase overfore, modo anexo. El formato lo pone la clase; aquí sólo va el contenido.
\documentclass[anexo]{../../comun/overfore}
\numsubdoc{1}
\informe{Informe Preparatorio N° 2}
\fecha{12 de octubre de 2026}
\externos{T14}
\begin{document}

\seccionsinnumero{Anexos del Subdocumento 1}
Estos anexos respaldan lo que declara el Subdocumento 1. Contienen las cartas de referencia de los mandantes de los cuatro proyectos declarados en el Formulario T-6, que se presenta en archivo propio, y el detalle del equipo clave, de la dotación de apoyo y de los especialistas externos. La Tabla \ref{tab:contenido-anexos} indica qué contiene cada anexo y qué sección lo cita.

\begin{tabla}{I{0.104}I{0.746}I{0.151}}{Contenido de los anexos.}{tab:contenido-anexos}
  {\cab{Anexo} & \cab{Contenido} & \cab{Sección que lo cita}}
  {Fuente: elaboración propia.}
A & Carta de referencia de Terminal Logístico Bahía Sur S.A. (NexoPuerto) & 1.4 \\
B & Carta de referencia de Alimentos Cordillera SpA (Planta Conecta) & 1.4 \\
C & Carta de referencia de Frío Austral Operaciones Ltda. (RedFrío 360) & 1.4 \\
D & Carta de referencia de Minera Pampa Tarel SpA (FaenaLink) & 1.4 \\
E & Equipo clave: roles, dedicación, período, responsabilidades y certificaciones & 1.4 y 1.5 \\
F & Dotación de apoyo y especialistas externos por perfil, cantidad y período & 1.2 y 1.5 \\
\end{tabla}
Los cuatro proyectos de referencia tienen carta de su mandante, y las cuatro cartas autorizan a la Comisión Evaluadora a verificar la información directamente con quien la firma. Los certificados institucionales y las cartas de los socios tecnológicos se acompañan en el Sobre N° 1, con los antecedentes técnicos habilitantes.

\anexo{Carta de referencia de Terminal Logístico Bahía Sur S.A.}
Santiago, 4 de septiembre de 2026

Señores\newline Comisión Evaluadora\newline Licitación N° TFEP-01/2026\newline Compañía Minera Altos de Aranda

Ref.: Carta de referencia del mandante para Overfore S.A.

De mi consideración:

Por medio de la presente, Terminal Logístico Bahía Sur S.A., en calidad de mandante, certifica que Overfore S.A. (RUT 77.482.961-K) ejecutó para nuestra empresa el proyecto NexoPuerto: trazabilidad y control de patio intermodal, entre 2023 y 2024, en el rubro de logística portuaria.

El proyecto consistió en una plataforma de trazabilidad de contenedores y carga a granel, con captura RFID, registro móvil fuera de cobertura, integración con nuestro ERP y con el sistema de despacho, tableros de operación por turno y auditoría de eventos. La arquitectura implementada es híbrida, con nube pública y un gateway local en el terminal. Overfore S.A. actuó como contratista principal, con responsabilidad sobre el diseño, la ejecución y la operación posterior.

La solución se encuentra finalizada y en operación. Se comprometió y se mantiene un nivel de servicio de disponibilidad de 99,7 \% mensual para las transacciones de patio. El volumen de operación soportado es de 1,8 millones de eventos al año, 170 usuarios y 42 equipos móviles.

Autorizamos a la Comisión Evaluadora a verificar esta información directamente con la persona que suscribe.

Atentamente,

Carolina Vega

Jefa de Transformación Operacional\newline Terminal Logístico Bahía Sur S.A.\newline carolina.vega@terminalbahiasur.cl\newline +56 9 6102 1840

\anexo{Carta de referencia de Alimentos Cordillera SpA}
\enlargethispage{4\baselineskip}
Santiago, 4 de septiembre de 2026

Señores\newline Comisión Evaluadora\newline Licitación N° TFEP-01/2026\newline Compañía Minera Altos de Aranda

Ref.: Carta de referencia del mandante para Overfore S.A.

De mi consideración:

Por medio de la presente, Alimentos Cordillera SpA, en calidad de mandante, certifica que Overfore S.A. (RUT 77.482.961-K) ejecutó para nuestra empresa el proyecto Planta Conecta: integración OT y gestión de mantenimiento, entre 2022 y 2023, en el rubro de manufactura de proceso.

El proyecto consistió en la lectura segura de variables de producción mediante OPC UA en modo solo lectura, con servidor de borde industrial, integración con mantenimiento, gestión de detenciones y visualización de indicadores para la supervisión. La arquitectura implementada es híbrida, con servidores industriales on-premise y servicios en nube. Overfore S.A. actuó como contratista principal, con responsabilidad sobre el diseño, la ejecución y la operación posterior.

La solución se encuentra finalizada y en operación. Se comprometió y se mantiene un nivel de servicio de disponibilidad de 99,6 \% mensual para captura y consulta operacional. El volumen de operación soportado es de 96 activos, 3.500 etiquetas de proceso, 520.000 eventos al año y 95 usuarios.

La Comisión Evaluadora podrá confirmar los antecedentes de Planta Conecta con Mauricio Rojas, gerente de operaciones de Alimentos Cordillera SpA, mediante los datos de contacto indicados en esta carta.

Atentamente,

Mauricio Rojas

Gerente de Operaciones\newline Alimentos Cordillera SpA\newline mauricio.rojas@alimentoscordillera.cl\newline +56 9 6127 3395

\anexo{Carta de referencia de Frío Austral Operaciones Ltda.}
\enlargethispage{4\baselineskip}
Santiago, 4 de septiembre de 2026

Señores\newline Comisión Evaluadora\newline Licitación N° TFEP-01/2026\newline Compañía Minera Altos de Aranda

Ref.: Carta de referencia del mandante para Overfore S.A.

De mi consideración:

Por medio de la presente, Frío Austral Operaciones Ltda., en calidad de mandante, certifica que Overfore S.A. (RUT 77.482.961-K) ejecutó para nuestra empresa el proyecto RedFrío 360: continuidad operacional para centros de distribución, entre 2021 y 2022, en el rubro de cadena de frío y distribución.

El proyecto consistió en una plataforma de telemetría, alarmas y trazabilidad de operaciones en centros de distribución, con almacenamiento local cifrado, sincronización automática, monitoreo y respaldos con recuperación ante desastres probada. La arquitectura implementada es híbrida, con nube pública, gateways de sitio y enlaces redundantes. Overfore S.A. actuó como contratista principal, con responsabilidad sobre el diseño, la ejecución y la operación posterior.

La solución se encuentra finalizada y en operación. Se comprometió y se mantiene un nivel de servicio de disponibilidad de 99,5 \% mensual para servicios críticos. El volumen de operación soportado es de 24 centros, 75 gateways, 680 dispositivos y 2,4 millones de lecturas al año.

Para cualquier consulta sobre la ejecución y operación de RedFrío 360, la Comisión Evaluadora puede contactar a Daniela Arancibia, subgerenta de tecnología de Frío Austral Operaciones Ltda., mediante el correo y teléfono señalados al cierre de esta carta.

Atentamente,

Daniela Arancibia

Subgerenta de Tecnología\newline Frío Austral Operaciones Ltda.\newline daniela.arancibia@frioaustral.cl\newline +56 9 6348 2751

\anexo{Carta de referencia de Minera Pampa Tarel SpA}
\enlargethispage{4\baselineskip}
Antofagasta, 4 de septiembre de 2026

Señores\newline Comisión Evaluadora\newline Licitación N° TFEP-01/2026\newline Compañía Minera Altos de Aranda

Ref.: Carta de referencia del mandante para Overfore S.A.

De mi consideración:

Por medio de la presente, Minera Pampa Tarel SpA, en calidad de mandante, certifica que Overfore S.A. (RUT 77.482.961-K) ejecutó para nuestra empresa el proyecto FaenaLink: registro de carguío y acreditación de contratistas en faena, entre 2024 y 2025, en el rubro de la minería del cobre a rajo abierto.

El proyecto consistió en una plataforma de registro del ciclo de carguío y transporte en terreno, con operación sin cobertura, acreditación de contratistas con control en portería, integración con el despacho de flota y con nuestro ERP, y tableros de operación por turno. La arquitectura implementada es híbrida, con nube pública y un servidor de borde en la faena. Overfore S.A. actuó como contratista principal, con responsabilidad sobre el diseño, la ejecución y la operación posterior.

El personal de Overfore S.A. trabajó en nuestra faena en operación, en régimen de turnos 7x7, acreditado como contratista conforme al Reglamento de Seguridad Minera y a nuestros procedimientos internos. El levantamiento funcional del proyecto estuvo a cargo de Patricio Figueroa Gallardo.

La solución se encuentra finalizada y en operación. Se comprometió y se mantiene un nivel de servicio de disponibilidad de 99,5 \% mensual. El volumen de operación soportado es de 46 camiones, 3.200 personas acreditadas, 1,4 millones de eventos al año y 260 usuarios.

La verificación de esta referencia puede coordinarse con Rodrigo Saavedra, Superintendente de Sistemas Mina, mediante los datos de contacto que se indican al cierre de la carta.

Atentamente,

Rodrigo Saavedra

Superintendente de Sistemas Mina\newline Minera Pampa Tarel SpA\newline rodrigo.saavedra@pampatarel.cl\newline +56 9 6215 4470

\anexo{Equipo clave}
El equipo clave reúne a nueve personas que ejercen once roles, desde la dirección del proyecto hasta la gerencia del servicio en la operación. La Tabla \ref{tab:roles-dedicacion-periodo-responsabilidad-equipo-clave} indica, para cada rol, quién lo ejerce, con qué dedicación y en qué período, y de qué responde en el proyecto.

\begin{tabla}{I{0.181}I{0.181}I{0.329}I{0.309}}{Roles, dedicación, período y responsabilidad del equipo clave.}{tab:roles-dedicacion-periodo-responsabilidad-equipo-clave}
  {\cab{Rol} & \cab{Integrante} & \cab{Dedicación y período} & \cab{Responsabilidad en el proyecto}}
  {Fuente: elaboración propia.}
Jefe de Proyecto & Claudio Cabello Vásquez & Exclusiva en la implementación, M1–M21. La reversibilidad y la devolución en la operación pasan al Gerente de Servicio. & Dirección general, relación formal con la contraparte y control de alcance y plazo \\
Arquitecto de Solución & Benjamín Rodríguez Pavez & Alta en el diseño, M1–M21; permanente en el Comité de Arquitectura, M1–M56 & Componentes lógicos, nube, red física, data centers y servidores de borde \\
Líder Funcional & Patricio Figueroa Gallardo & Completa en la implementación, M1–M21 & Contexto del problema, requerimientos, alcance, exclusiones y trazabilidad \\
Líder de Datos & Benjamín Velásquez Farías & Permanente en la implementación, M1–M21; en la operación, sólo los informes de las innovaciones 1 y 5, hasta M54 y M46 respectivamente & Modelo y gobierno de datos, migración y trazabilidad del mineral \\
Líder de Implantación y Gestión del Cambio & Benjamín Velásquez Farías & Desde M2, con la dedicación exigida desde M8; en la operación, la capacitación continua hasta M56 & Puesta en marcha, marchas blancas y adopción en faena \\
Líder de Integración & Marcelo Flores Coloane & Permanente en la implementación, M1–M21 & Bus de eventos y adaptadores con ERP, despacho, laboratorio y telemetría \\
Líder de Desarrollo & Gabriel Sanzana Cuibin & Completa en la implementación, M1–M21; mantención en la operación & Servicios de negocio, aplicación de captura desconectada y flujos del sistema \\
Encargado de Seguridad de la Información & David Martínez Canto & Permanente, M1–M56 & Arquitectura de seguridad, confianza cero, cifrado y segregación TI/TO \\
Líder de Operación / SRE & Lucas Zamora Gatica & Desde M6; permanente en la operación, M21–M56 & Mesa de ayuda y centros de operaciones de red y de seguridad 24x7; continuidad 24x7, observabilidad, recuperación ante desastres y niveles de servicio \\
Gerente de Servicio & Claudio Cabello Vásquez & Dedicado durante la operación, M21–M56 & Contraparte permanente del cliente; gerencia el servicio y responde por la reversibilidad y la devolución desde el mes 21 \\
Líder de Calidad y Pruebas & Kavon Kermani Órdenes & Permanente, M1–M56 & Aseguramiento de calidad, planes de prueba y validación previa al paso a producción \\
\end{tabla}
Claudio Cabello Vásquez ejerce dos roles, Jefe de Proyecto y Gerente de Servicio, sólo desde la primera ola del paso de la Etapa 2 hasta el acta de H12. Benjamín Velásquez Farías ejerce dos roles, Líder de Datos y Líder de Implantación, con las dedicaciones de la Tabla \refext{T14}{tab:roles-equipo-clave-dedicacion-paquetes-cargo} del Formulario T-14.

Cada integrante tiene certificaciones vigentes ligadas a su rol en Aranda. La Tabla \ref{tab:certificaciones-competencias-equipo-clave} las presenta con el trabajo del proyecto en que se aplican.

\begin{tabla}{I{0.186}I{0.186}I{0.309}I{0.318}}{Certificaciones y competencias del equipo clave.}{tab:certificaciones-competencias-equipo-clave}
  {\cab{Integrante} & \cab{Rol} & \cab{Certificación o competencia} & \cab{Aplicación en Aranda}}
  {Fuente: elaboración propia.}
Claudio Cabello Vásquez & Jefe de Proyecto & PMP, ITIL 4 Foundation, COBIT 2019 Foundation, Scrum Master & Dirección del proyecto, control de alcance y de plazo, y relación formal con la contraparte \\
Benjamín Rodríguez Pavez & Arquitecto de Solución & Azure Solutions Architect Expert, TOGAF Foundation, ITIL 4 Foundation, Certified Data Centre Professional & Arquitectura híbrida, topología de nube y de borde, data centers y servidores de faena \\
Patricio Figueroa Gallardo & Líder Funcional & PMP, IREB CPRE Foundation, ISTQB Foundation; levantamiento funcional de FaenaLink, minería del cobre, 2024–2025 & Levantamiento del problema, requerimientos, alcance y criterios de aceptación \\
Benjamín Velásquez Farías & Líder de Datos e Implantación y Gestión del Cambio & Azure Data Engineer Associate, EDB Certified Associate PostgreSQL & Modelo y gobierno de datos, trazabilidad del mineral y puesta en marcha en faena \\
Marcelo Flores Coloane & Líder de Integración & OPC UA Fundamentals, ISA/IEC 62443 Fundamentals, Certified Ethical Hacker (CEH) & Bus de eventos y adaptadores hacia ERP, despacho, laboratorio y telemetría \\
Gabriel Sanzana Cuibin & Líder de Desarrollo & Azure Solutions Architect Expert, Azure DevOps Engineer Expert, EDB Certified Associate PostgreSQL, Scrum Master & Servicios de negocio, captura desconectada en terreno y flujos del sistema \\
David Martínez Canto & Encargado de Seguridad de la Información & CISSP, COBIT 2019 Foundation, ITIL 4 Foundation, Fortinet FCP & Arquitectura de seguridad, confianza cero, cifrado y segregación TI/TO \\
Lucas Zamora Gatica & Líder de Operación y SRE & ITIL 4 Foundation, Azure Solutions Architect Expert, Linux RHCSA, Cisco CCNA & Continuidad operacional, observabilidad, recuperación ante desastres y niveles de servicio \\
Kavon Kermani Órdenes & Líder de Calidad y Pruebas & ISTQB Foundation, ITIL 4 Foundation, auditor interno ISO 9001 & Aseguramiento de calidad, planes de prueba y validación previa al paso a producción \\
\end{tabla}
El equipo clave suma dos certificaciones en gestión de proyectos, cinco en ITIL 4, dos en COBIT, tres de Azure Solutions Architect Expert, dos en seguridad de la información, dos en PostgreSQL, que es el motor de base de datos de la solución, y dos en pruebas de software. Cada certificación se acredita con copia del certificado vigente y el código de verificación del organismo emisor.

\anexo{Dotación de apoyo y especialistas externos}
La dotación de apoyo es personal propio de Overfore, asignado por perfil, cantidad y período. Proviene del área de desarrollo, del centro de operaciones y de la base de Antofagasta, y se amplía con contrataciones programadas antes de cada período. La Tabla \ref{tab:dotacion-apoyo-overfore-asignada-proyecto} la detalla.

\begin{tabla}{I{0.171}I{0.175}I{0.273}I{0.165}I{0.216}}{Dotación de apoyo de Overfore asignada al proyecto.}{tab:dotacion-apoyo-overfore-asignada-proyecto}
  {\cab{Perfil} & \cab{Personas} & \cab{Período} & \cab{Base} & \cab{Trabajo que ejecuta}}
  {Fuente: elaboración propia.}
Desarrollador & 8 en M3–M11, 2 en M12, 7 en M13–M17 y 2 en M18–M20. Período: M3–M20. & M3–M20; apoyo inicial a ambientes y flujo de integración en M1–M2 & Casa matriz, Santiago & Módulos, conectores, canales e innovaciones, en los tres equipos Scrum \\
Desarrollador & 1 & M21–M56 & Casa matriz, Santiago & Mantención correctiva, evolutiva y normativa \\
Analista de pruebas & 1 & M1–M20; y en la operación una parte de su jornada para las puertas de calidad de cada versión (1.8.8.4) & Casa matriz, Santiago & Pruebas automatizadas, regresión y aceptación \\
Analista de datos & 1 en M2–M12 y M18–M20, 2 en M13–M17 & M2–M20; en la operación, una parte de su jornada para la huella de carbono por lote (1.8.7.3) y las evaluaciones de las innovaciones hasta M54 & Casa matriz, Santiago & Perfilado, migración, tableros e indicadores \\
Digitalizador & 4 & M10–M15; rotando en tres puestos de escaneo & Faena & Digitalización de los formularios de seguridad, con contrato a plazo fijo \\
Técnico de terreno & 1 en M6–M12; 8 en M13–M15; 3 en M16–M18; 10 en M19–M20; 10 en M21–M23; 5 desde M24 & M6–M12 por campaña; M13–M56 continuo, en turnos 7x7 & Antofagasta & Instalación de dispositivos y soporte presencial en faena, en turnos 7x7 \\
Analista de operación & 8 en las semanas 6 a 13 de la marcha blanca (M14–M15); 2 en M16–M17; 3 en M18; 10 en M19–M20; 10 en M21–M23; 10 desde M24 & M13–M23 en faena (acompañamiento de las marchas blancas, estabilización y soporte de la Etapa 1); M21–M56 en turnos 24x7 & Centro de operaciones, Santiago & Mesa de ayuda y NOC 24x7 \\
Analista de seguridad & 5 & M21–M56 en turnos; participación previa desde M2 en amenazas y vulnerabilidades & Centro de operaciones, Santiago & SOC 24x7 \\
\end{tabla}
La dotación de apoyo de desarrollo, pruebas, datos y digitalización llega a catorce personas como máximo entre los meses 10 y 15. Con los analistas de operación y los técnicos de la marcha blanca, la implementación llega a treinta personas en los meses 13 a 15. En la operación llega a veintiuna personas, cuando los servicios 24x7 exigen turnos continuos. A las catorce personas de desarrollo, pruebas, datos y digitalización se suma la participación parcial de los analistas de seguridad y de operación. Esa participación cubre las amenazas, los libros de operación y el paso a producción.

La dotación de la operación se calcula por puesto de servicio. Un puesto 24x7 cubre dos turnos de 12 horas por día, es decir 8.760 horas al año, y una persona con jornada de 40 horas semanales, vacaciones y feriados trabaja unas 1.824 horas al año. Cada puesto necesita por eso cinco personas, que también cubren los relevos. La mesa de ayuda y el NOC ocupan un puesto cada uno, con los diez analistas de operación, y el SOC ocupa un puesto con los cinco analistas de seguridad. En la faena hay un técnico en el turno de día y otro en el de noche, con jornadas de 12 horas y turnos 7x7. El puesto 24x7 necesita cinco técnicos, que cubren también vacaciones, licencias y campañas de mantención durante los 36 meses de operación.

Los especialistas externos cubren trabajos puntuales para los que no conviene mantener personal permanente. La Tabla \ref{tab:especialistas-externos} indica el trabajo de cada perfil, su etapa y su forma de contratación.

\begin{tabla}{I{0.208}I{0.301}I{0.192}I{0.298}}{Especialistas externos.}{tab:especialistas-externos}
  {\cab{Perfil} & \cab{Trabajo que ejecuta} & \cab{Etapa} & \cab{Forma de contratación}}
  {Fuente: elaboración propia.}
Auditor de intrusión independiente & Pruebas de intrusión de la implementación y campañas anuales & Implementación y operación & Empresa sin relación con el desarrollo, a precio fijo por campaña \\
Diseñador UX/UI & Investigación con usuarios, sistema de diseño, prototipos y pruebas de usabilidad & Implementación & Por horas \\
Científico de datos & Clasificador asistido de detenciones y su verificación & Implementación y operación hasta M33 & Por horas \\
Relator de capacitación & Capacitación por perfil y turno, y jornadas semestrales de actualización & Implementación y operación & Por horas \\
Consultor en gestión del cambio & Diagnóstico de impacto, red de agentes de cambio y procedimientos rediseñados & Implementación & Por horas \\
Redactor técnico & Manuales, guías rápidas y base de conocimiento & Implementación & Por horas \\
Abogado en datos personales y derecho laboral & Acuerdo con el sindicato, contratos, evaluación de impacto y adecuaciones normativas & Implementación y operación & Por horas \\
Analista de control de proyectos & Plan para la dirección, informes de valor ganado y actas de comité & Implementación & Por horas \\
Especialista en continuidad del negocio & Plan de continuidad y sus actualizaciones anuales & Implementación y operación & Por horas \\
Especialista en pruebas de rendimiento & Pruebas de carga, estrés y resistencia & Implementación y operación & Por horas \\
Especialista en radiofrecuencia & Estudio de cobertura de la red del rajo y su revisión anual & Implementación y operación & RedNorte Industrial SpA, por horas \\
Montajista eléctrico y de comunicaciones & Montaje de la sala técnica, los gabinetes y los enlaces & Implementación & RedNorte Industrial SpA, por horas \\
\end{tabla}
Los especialistas externos aportan conocimiento acotado y no reemplazan a ningún rol del equipo clave. Toda contratación externa requiere autorización previa y escrita del cliente.

\end{document}
